% Generated draft: requires mathematical and source review.
\documentclass[11pt]{article}
\usepackage[margin=25mm]{geometry}
\usepackage{fontspec}
\setmainfont{Latin Modern Roman}
\usepackage{amsmath,amssymb,mathtools,unicode-math}
\setmathfont{Latin Modern Math}
\usepackage{enumitem,needspace,xurl,hyperref}
\setlength{\parindent}{0pt}
\setlength{\parskip}{5pt}
\setlength{\emergencystretch}{3em}
\newcommand{\N}{\mathbb{N}}
\newcommand{\Z}{\mathbb{Z}}
\newcommand{\Q}{\mathbb{Q}}
\newcommand{\R}{\mathbb{R}}
\newcommand{\C}{\mathbb{C}}
\newcommand{\F}{\mathbb{F}}
\newcommand{\A}{\mathbb{A}}
\newcommand{\PP}{\mathbb{P}}
\newcommand{\E}{\mathbb{E}}
\newcommand{\eps}{\varepsilon}
\newcommand{\dd}{\,\mathrm{d}}
\begin{document}
\section*{UnsolvedMath ID: 30001241}
\textbf{Source problem code:} catalogue entry
\subsection{Definitions and mathematical statement}
Let \(a,b,c\in\Z_{>0}\) satisfy \(\gcd(a,b)=\gcd(a,c)=\gcd(b,c)=1\). To specify the algebraic notation, fix a field \(k\) and put \(A=k[x,y,z]\). The choice and characteristic of \(k\) are not specified in the supplied source record; the following construction is therefore relative to this auxiliary choice. Define
\[
\varphi_{a,b,c}:A\longrightarrow k[t],\qquad x\longmapsto t^a,\quad y\longmapsto t^b,\quad z\longmapsto t^c,
\qquad P=P(a,b,c):=\ker(\varphi_{a,b,c}).
\]
Thus \(P\) is the prime ideal defining the affine monomial curve parametrized by \((t^a,t^b,t^c)\). Set \(P^{(0)}=A\), and for every integer \(m\geq1\) define the symbolic power
\[
P^{(m)}:=P^mA_P\cap A.
\]
With \(T\) an indeterminate recording the grading, its symbolic Rees ring is
\[
R_s(P):=\bigoplus_{m\geq0}P^{(m)}T^m\subseteq A[T].
\]
\textbf{Problem.} Characterize the pairwise coprime positive triples \((a,b,c)\) for which \(R_s(P(a,b,c))\) is finitely generated as an \(A\)-algebra, equivalently Noetherian, with the necessary hypotheses on \(k\) made explicit. Determine the relationship between this finite-generation problem and Nagata's conjecture for general points in \(\mathbb P^2\).

\textbf{Unresolved qualifications.} The source does not specify the ground field, its characteristic, or whether the intended characterization is uniform over fields. It also does not state a precise version of Nagata's conjecture, the number or multiplicities of the points, the meaning of generality, or a proposed implication or equivalence. Consequently, the requested relationship cannot be expressed as a uniquely determined logical assertion from the supplied record alone. The positive-integer hypotheses permit an exponent equal to \(1\); no exclusion of such cases is supplied.
\subsection{Short English statement}
For pairwise coprime positive integers \(a,b,c\), characterize when the symbolic Rees ring of the monomial curve \((t^a,t^b,t^c)\) is Noetherian, and explain its relationship to Nagata's conjecture for general plane points. The ground-field assumptions and the precise conjectural relationship remain unspecified in the supplied record.
\subsection{Sources}
\begin{itemize}
\item Supplied local UnsolvedMath registry record, permanent ID \(30001241\), exact source problem code \texttt{catalogue entry}, titled ``Symbolic Rees Rings and Nagata's Conjecture,'' in the Oberwolfach Reports Open Problems collection. The record assigns the year 2009 and supplies no named proposer.
\item \emph{Kommutative Algebra} (2009), Oberwolfach Reports, Report 22, pp.\ 1157--1224, DOI \texttt{10.4171/owr/2009/22}. Supplied official report reference: \url{https://doi.org/10.4171/owr/2009/22}.
\item Supplied literature reference labeled ``Kurano-Matsuoka on space monomial curves'': \url{https://arxiv.org/abs/0801.3896}.
\item Supplied literature reference labeled ``Kurano-Nishida on infinitely generated examples'': \url{https://arxiv.org/abs/1705.09865}.
\item Literature assessment dated exactly \texttt{2026-08-21}: the supplied curation labels the problem open and describes its connection to Nagata-type negative-curve questions, with later papers resolving families and examples but not the general problem. This is supplied dated curation, not an independently verified literature assessment; its time-sensitive status and cited claims require independent verification.
\item The registry labels the statement \texttt{corrected\_verified} and reports that its mathematical content was checked against the official Oberwolfach report, with notation, context, and question form normalized without changing the source-level claim. This verification description is supplied metadata and has not been independently confirmed for this draft.
\end{itemize}
\end{document}
